\section{伪模}

记号约定:$A$是伪环,$A^{\prime}$是$A$的幺化.

\begin{definition}{左$A$-伪模,右$A$-伪模}{}
	设$E$是交换群(加法记号).
	\begin{enumerate}
		\item 设$A$的加法群左作用于$E$.若该作用只是半群作用,对$E$的加法分配,对$A$的加法分配,则称配备该左作用的$E$为左$A$-模.
		\item 设$A$的加法群左作用于$E$.若该作用只是半群作用,对$E$的加法分配,对$A$的加法分配,则称配备该右作用的$E$为右$A$-模.
	\end{enumerate}
\end{definition}

\begin{proposition}{伪模的幺化}{unital-of-pseudomodules}
	设$E$是左$A$-伪模.定义$A^{\prime}$在$E$上的左作用为$\left(\left(n,a\right),x\right)\mapsto nx+ax$,则$E$是左$A^{\prime}$-模,并且该$A^{\prime}$-模唯一.
\end{proposition}

\begin{definition}{伪模的幺化}{}
	设$E$是左$A$-伪模.我们称命题(\ref{prop:unital-of-pseudomodules})中得到的$A^{\prime}$-模是$E$的幺化.
\end{definition}





















